\Qhead{1.10}{Распределённые автоколебательные системы. Лазер}{Распределенные автоколебательные системы. Лазер как пример такой системы. Условия самовозбуждения. Одномодовый и многомодовый режимы генерации.}
\Subfile{1.10.1}{sections/01_oscillations/q10/s01}
\Subfile{1.10.2}{sections/01_oscillations/q10/s02}

\recap{Моды распределённого резонатора получают энергию от активной среды. Генерация начинается при превышении усиления над потерями; насыщение ограничивает амплитуду. Двухмодовый фазовый портрет показывает вытеснение, устойчивое сосуществование или бистабильность.}
\board{Резонатор с накачкой; модальное уравнение и порог; две прямые нулевого роста и три фазовых портрета.}
\oral{Почему слабая конкуренция ещё не гарантирует существования двухмодового режима?}
\bait{Пространственное выжигание инверсии: разные стоячие волны расходуют усиление в разных местах.}
\sources{Балакший~\cite{balakshy}, PDF-с.~25--27, листы~23--25: He--Ne-лазер, разложение поля и поляризации, укороченные уравнения и фазовые портреты. PDF-с.~28, лист~26: отличие многомодовости от синхронизации. Условие обхода сверено со слайдом «Условие самовозбуждения лазера»~\cite{optoelectronics}; насыщение --- с Kärtner~\cite{kartner}, §4.2, печатная с.~132. Обозначения и проверка устойчивости приведены к единой нормировке.}
